\section{Appendix}
\label{sec:appendix}

\subsection{Audit-first disclosure for missingness and edge cases}

For all meta-label and tag-derived quantities, NA denotes a system/format missing value (e.g., field not
transmitted, corrupted record). Our protocol treats NA as an auditable disclosure item: NA is reported as
a rate with explicit denominators, and affected tasks are not silently removed from analysis.

Separately from NA, we disclose \textbf{non-compliant} multi-select meta-label records as Type~4
process/format failures (e.g., empty selection when ``at least one'' is required, or mutually-exclusive
tags co-selected such as \texttt{trivial}+\{other difficulty\} / \texttt{acceptable}+\{other failure modes\}).
These events are reported with explicit denominators and example tasks to support audit and downstream
counter-example curation.

When UI-level hard validation is enabled, many Type~4 events are prevented \emph{at submission time}.
To avoid ``hiding'' these errors, we additionally report the number of blocked submission attempts and the
distribution of rejection reasons as process evidence.

For \texttt{PreScreen\_semi}, we additionally disclose the source composition of misleading initializations.
Synthetic perturbation operators constitute the primary source. A small, pre-registered natural-failure trap
bank may be used as a supplement on a disjoint image pool; the same \texttt{base\_task\_id} is never reused
across \texttt{PreScreen\_manual} and \texttt{PreScreen\_semi}. If natural failures are insufficient for a given
failure family, the realized quota for that family is reduced and the remainder is filled by the corresponding
synthetic family. Planned and realized source quotas are both disclosed.

For OOS handling, we distinguish clearly judgeable scope-gate/stress items from high-ambiguity cases. Tasks that
are clearly judgeable as OOS or insufficient evidence but lack a stable geometric reference may enter a fixed small
manual scope-gate subset or a smaller semi stress subset, both disclosed separately from the main geometric trap
pool. Highly ambiguous OOS cases are retained only in an ambiguity audit bank.

For the weighted consensus operator (Algorithm~\ref{alg:weighted-tag-consensus}), the event
$W_{\mathrm{tot}}=0$ is also disclosed. This paper fixes the corresponding policy prior to the Main
experiment and reports its frequency. The main analysis records this event as NA rather than silently falling
back to an unweighted vote; sensitivity checks may additionally compare a ``fallback-to-unweighted''
variant, but this comparison is reported only as a robustness disclosure and is not used to select results.

For the lock-in of reliability weighting, we further disclose the anchor design used by
\texttt{PreScreen\_manual}: a 30-image pool containing 20--22 expert-referenced manual anchors and 8--10
random non-anchor images. The main cross-fitting rule is repeated balanced two-fold splitting on the manual
anchors; three-fold stratified splitting is reported only as a sensitivity analysis, while five-fold CV is not
used as a main lock-in protocol under the current anchor budget.

\subsection{Pre-registered sensitivity grids (reporting only)}

When reporting sensitivity analyses (e.g., varying the consensus threshold $\tau_c$ or the cap of a weight
mapping), we follow a pre-registered grid and report the full grid (or a pre-registered summary such as
median and interquartile range). No post-hoc choice is made based on observed improvements.
